\manualchapter{Patches and troubleshooting}{sec:operation}
\subsection{Layered tone and noise}
Use Triangle on voice 1 and Sample+Hold on voice 2. Give them the same pitch
via normalling, then tune voice 2 higher. Patch separate envelopes to the two
Level inputs: a short noise envelope makes a percussive attack over a longer
tonally stable voice. Take each output separately for external mixing, or
use the forward sum and reduce both levels to leave headroom.

The four selected shapes are saved alongside parameters in patches and
presets. Reset restores Sine on every oscillator; randomize chooses each voice's shape
as well as changing parameters. All polyphonic channels of a voice share its
shape, while their pitch and amplitude can differ.

If a voice is silent, check its Level CV and output routing. If the sound
becomes noisy after a shape change, choose Sine explicitly before checking
tuning. The sample-and-hold and shift-register modes are not sine oscillators
with added distortion; their clocked patterns are the intended source.

\subsection{Cable channels, outputs and saved patches}
The widest input selects 1--16 polyphonic chip instances. Voice columns are
oscillators within each instance, not separate cable channels. Inputs read
the matching channel; supply matching pitch and level channel counts rather
than expecting a mono envelope to repeat across a polyphonic pitch cable.

Each output adds the preceding output when that earlier socket is unpatched.
The accumulated signal clips at approximately $\pm5$\,V. To hear the complete
mix, use the last output and leave earlier outputs empty. Patching an earlier
output removes that signal from the following mix. Lower source levels if
the sum clips.

Rack saves the panel parameters. Reset initializes the chip; oscillator
phase is not preserved by a patch save. Pitch, FM and level are read each sample; the oscillators render their
waveforms in internal blocks.
